\section{Model, boundary-value problem, and certified solution}
\label{sec:model}

We work on a prescribed four-dimensional de Sitter static patch, with metric
signature \((+---)\),
\[
  ds^2=N(r)\,dt^2-N(r)^{-1}\,dr^2-r^2d\Omega_2^2,
  \qquad
  N(r)=1-\frac{r^2}{R_{\mathrm{dS}}^2}.
\]
The matter field is \(U\colon\mathcal M\to SU(2)\), with
\(R_\mu=U^\dagger\nabla_\mu U\), and our massive Skyrme convention is
\[
  S_{\mathrm{Sk}}
  =\int_{\mathcal M}d^4x\,\sqrt{-g}\left[
    -\frac{f_\pi^2}{16}\operatorname{Tr}(R_\mu R^\mu)
    +\frac{1}{32e^2}\operatorname{Tr}
       \bigl([R_\mu,R_\nu][R^\mu,R^\nu]\bigr)
    +\frac{f_\pi^2m_\pi^2}{16}
       \operatorname{Tr}(U+U^\dagger-2)
  \right].
\]
In particular, \(m_\pi\) is the mass of the linearized pion field; the
coefficient of the last term is not the convention in which this mass is
\(\sqrt{2}m_\pi\).

For the static hedgehog
\[
  U(r,\widehat x)=\cos F(r)
  +i\,\boldsymbol\tau\mathbin{\cdot}\widehat x\,\sin F(r),
\]
introduce
\[
  x=ef_\pi r,
  \qquad
  \mu=\frac{m_\pi}{ef_\pi},
  \qquad
  \lambda=\frac{1}{(ef_\pi R_{\mathrm{dS}})^2},
  \qquad
  N(x)=1-\lambda x^2,
\]
and abbreviate \(s=\sin F\) and \(u=x^2+8s^2\).  The bulk static mass is
\[
  M_{\mathrm{bulk}}=\frac{4\pi f_\pi}{e}
     \int_0^{x_w}\mathcal E(x)\,dx,
\]
where
\[
  \mathcal E
  =\frac18Nu(F')^2+\frac14\sin^2F
    +\frac{\sin^4F}{2x^2}
    +\frac{\mu^2x^2}{4}(1-\cos F).
\]
Here and below a prime denotes \(d/dx\).  Its Euler--Lagrange equation is
\begin{equation}
  \bigl(NuF'\bigr)'
  =\left(4N(F')^2+1+\frac{4\sin^2F}{x^2}\right)\sin(2F)
    +\mu^2x^2\sin F.
  \label{eq:model-profile}
\end{equation}
The apparent singularity at \(x=0\) is removable on the regular hedgehog
branch described below.

\subsection*{The prescribed Dirichlet boundary}

We truncate the static domain at \(x=x_w\), strictly inside the Killing horizon
\(x_c=\lambda^{-1/2}\), and impose the Dirichlet constraint
\(U|_{x_w}=\mathbf 1_2\).  In the unit-charge hedgehog sector the
boundary data are therefore
\begin{equation}
  F(0)=\pi,
  \qquad
  F(x_w)=0.
  \label{eq:model-boundary}
\end{equation}
No wall equation of motion is imposed in this paper:
\eqref{eq:model-boundary} is prescribed data,
not a claim that a particular finite-tension membrane dynamically realizes the
boundary.

For the certified problem,
\begin{equation}
  \mu^2=1,
  \qquad
  \lambda=\frac1{400},
  \qquad
  x_w=4,
  \qquad
  x_c=20.
  \label{eq:model-parameters}
\end{equation}
Thus the complete boundary-value problem is
\eqref{eq:model-profile}--\eqref{eq:model-parameters} on \([0,4]\).
This parameter point is a validation benchmark, not a phenomenological fit.
The choice \(\mu=1\) keeps the pion and Skyrme length scales comparable, while
\(\sqrt\lambda x_w=1/5\) gives the rational wall lapse
\(N(x_w)=24/25\) and keeps the complete domain well inside the horizon.  It
therefore tests nonzero mass, curvature, and finite confinement without a
near-horizon singular limit.  The paper does not infer a distinct de Sitter
effect from this single point.  A qualitative local continuation family is
proved in \cref{cor:local-parameter-branch}, without an explicit box width.

\subsection*{Regularity class and the augmented chart}

We record the function spaces because they are part of the scope of the
uniqueness statement.  Set
\[
  a=\frac1{16},
  \qquad
  I=[a,4].
\]
The removable origin is treated in the variables
\[
  t=x^2,
  \qquad
  \pi-F=xu(t),
  \qquad
  p_b(t)=\frac{d(\pi-F)}{dx}.
\]
For a positive shooting slope \(b=-F'(0)\), the validated Volterra chart has
\begin{equation}
  p_b(t)=b-3c(b)t+t^2\rho_b(t),
  \qquad
  \|\rho_b\|_{L^\infty(0,a^2)}\le 20.
  \label{eq:model-origin-ball}
\end{equation}
The cubic coefficient appearing here is exactly
\begin{equation}
  c(b)=\frac{b\left[
       \mu^2-4\lambda+\left(\frac43-24\lambda\right)b^2
       +\frac83b^4\right]}
       {10(1+8b^2)}.
  \label{eq:model-cubic}
\end{equation}
Consequently
\begin{equation}
  \left|F(x)-\bigl(\pi-bx+c(b)x^3\bigr)\right|
     \le 4x^5,
  \qquad
  \left|F'(x)-\bigl(-b+3c(b)x^2\bigr)\right|
     \le 20x^4.
  \label{eq:model-origin-bounds}
\end{equation}
These bounds hold on \([0,a]\).  The Volterra contraction selects one regular
solution inside
this weighted ball for each certified \(b\); it does not assert uniqueness
among all conceivable singular-origin function classes.

Let \(F_b^{\mathrm{org}}\) denote that regular origin solution, and define its
cutoff data
\[
  \Phi(b)=F_b^{\mathrm{org}}(a),
  \qquad
  \Gamma(b)=(F_b^{\mathrm{org}})'(a).
\]
On \(I\), the augmented formulation consists of
\eqref{eq:model-profile} together with
\begin{equation}
  F(a)=\Phi(b),
  \qquad
  F(4)=0,
  \qquad
  F'(a)=\Gamma(b).
  \label{eq:model-augmented-boundary}
\end{equation}
On the positive-radius interval the Dirichlet correction space is
\[
  X_D=\{h\in W^{2,\infty}(I):h(a)=h(4)=0\}.
\]
The exact center slope used by the certificate is
\[
  \overline b=\frac{1462763601523}{925827031517}.
\]
Let \(\widehat F\) be the archived globally \(C^2\), piecewise-rational
polynomial profile, and set
\[
  \chi_a(x)=\frac{4-x}{4-a},
  \qquad
  \chi_4(x)=\frac{x-a}{4-a},
\]
\begin{equation}
  \overline u
    =\widehat F-\chi_a\widehat F(a)-\chi_4\widehat F(4),
  \qquad
  \overline F=\overline u+\chi_a\Phi(\overline b).
  \label{eq:model-center}
\end{equation}
Thus \(\overline F(a)=\Phi(\overline b)\) and
\(\overline F(4)=0\) exactly; neither endpoint residual of the generated
spline is discarded.  The inverse used by the certificate is built instead
from the Jacobi operator at the raw rational spline \(\widehat F\).  Put
\[
  \mathsf p(x,F)=N(x)\bigl(x^2+8\sin^2F\bigr),
  \qquad
  W(x,F)=\sin^2F+\frac{2\sin^4F}{x^2}
          +\mu^2x^2(1-\cos F),
\]
and let
\begin{equation}
  A_0h=-(P_0h')'+Q_0h,
  \qquad
  P_0=\mathsf p(x,\widehat F),
  \qquad
  Q_0=W_{FF}(x,\widehat F)
       +\frac12\mathsf p_{FF}(x,\widehat F)(\widehat F')^2
       -(\mathsf p_F(x,\widehat F)\widehat F')'.
  \label{eq:model-jacobi}
\end{equation}
The certificate proves that the graph norm
\[
  \|h\|_{A_0}=\|A_0h\|_{L^\infty(I)}
\]
is equivalent to the usual \(W^{2,\infty}\) norm on \(X_D\), and hence
controls \(h\) and \(h'\) pointwise.  The endpoint-corrected
\(\overline F\) remains the center of the nonlinear chart below.  The
certificate separately encloses the coefficient change from
\(D\mathcal G(\widehat F)=A_0\) to
\(D\mathcal G(\overline F)\); this is the principal contribution to the
reported derivative-mismatch bound \(Z_0\).

For completeness, the local neighborhood used by the theorem can be stated
without suppressing the shooting variable.  Let \(H\) be the certified
homogeneous Jacobi solution
\[
  A_0H=0,
  \qquad
  H(a)=\Phi'(\overline b),
  \qquad
  H(4)=0.
\]
In local coordinates
\((w,\beta)\in X_D\times\mathbb R\), the physical profile and shooting slope
are
\begin{equation}
  \begin{split}
  F_{w,\beta}
    ={}&\overline F+w+\beta H\\
      &+\chi_a\bigl[\Phi(\overline b+\beta)-\Phi(\overline b)
                    -\Phi'(\overline b)\beta\bigr],
  \qquad
  b=\overline b+\beta .
  \end{split}
  \label{eq:model-chart}
\end{equation}
This chart enforces the two value conditions in
\eqref{eq:model-augmented-boundary} for every coordinate pair.  The derivative
condition remains the scalar augmented equation; at a zero it gives the
\(C^1\) join to the regular origin family.  With \(\omega=3/4\), the certified
closed ball is
\begin{equation}
  \mathcal B_r
  =\left\{(w,\beta):
      \max\bigl(\|w\|_{A_0},\omega|\beta|\bigr)\le r\right\},
  \qquad
  r=\frac1{250}.
  \label{eq:model-ball}
\end{equation}
In particular, its slope interval is exactly
\begin{equation}
  \overline b+\left[-\frac{2}{375},\frac{2}{375}\right]
  =\left[
    \frac{546684696508091}{347185136818875},
    \frac{550388004634159}{347185136818875}
  \right].
  \label{eq:model-slope-interval}
\end{equation}

\begin{theorem}[Certified massive Dirichlet profile]
\label{thm:main}
For the fixed-background boundary-value problem
\eqref{eq:model-profile}--\eqref{eq:model-parameters}, there exists a pair
\((w_*,\beta_*)\in\mathcal B_{1/250}\) such that \(F_*=F_{w_*,\beta_*}\)
solves \eqref{eq:model-profile} and all three conditions in
\eqref{eq:model-augmented-boundary}.  Adjoining the certified regular-origin
piece gives a profile on \([0,4]\) obeying \eqref{eq:model-boundary}.  This pair
is the only solution represented in \(\mathcal B_{1/250}\).

The exact-rational Newton certificate, using \(\omega=3/4\), twelve-term
Taylor--Lagrange trigonometric enclosures, and outward rounding on the grid
with denominator \(10^{18}\), bounds the Newton defect \(Y\), the
approximate-inverse derivative mismatch \(Z_0\), and the Schur-composed
nonlinear Hessian bound \(Z_2\) by
\begin{equation}
  Y\le\frac{287}{125000},
  \qquad
  Z_0\le\frac{39}{500000},
  \qquad
  Z_2\le\frac{78733}{1000}.
  \label{eq:model-newton-bounds}
\end{equation}
For \(r=1/250\), the corresponding radii polynomial and contraction bound
satisfy
\begin{equation}
  \mathfrak p(r)
  =Y+(Z_0-1)r+\frac12Z_2r^2
    <-\frac{107}{100000},
  \qquad
  \mathfrak q(r)=Z_0+Z_2r<\frac{31501}{100000}<1.
  \label{eq:model-radii-closure}
\end{equation}
Therefore the self-map and contraction inequalities both close.

Moreover, \(F_*\) is strictly decreasing on \([0,4]\).  Its wall slope lies
in the exact archived rational enclosure
\begin{equation}
  F_*'(4)\in
  \left[
    -\frac{94644972770042989}{10^{18}},
    -\frac{43733242881735721}{5\cdot10^{17}}
  \right]
  \subset
  \left[-\frac{1893}{20000},-\frac{4373}{50000}\right].
  \label{eq:model-wall-slope}
\end{equation}
In particular, the wall slope is nonzero and negative.  Finally, define the
dimensionless interior rigid-rotor inertia by
\begin{equation}
  \mathcal I[F]
  =\frac{2\pi}{3}\int_0^4x^2\sin^2F
     \left[
       \frac1N+4(F')^2+\frac{4\sin^2F}{Nx^2}
     \right]dx.
  \label{eq:model-inertia}
\end{equation}
Directed interval quadrature on the same Newton tube, together with the
regular-origin estimate, proves the exact conservative enclosure
\begin{equation}
  \frac{21149}{1000}
  \le \mathcal I[F_*]
  \le \frac{48921}{1000}.
  \label{eq:model-inertia-bound}
\end{equation}
Thus the matter contribution to the physical collective inertia,
\(I_{\mathrm{bulk}}=\mathcal I[F_*]/(e^3f_\pi)\), is finite and strictly
positive.
\end{theorem}

The paper uses the source-consistent combined AU.1/AU.2 archive
\path{experiments/skyrmion_au2_global_tail_exact_certificate.json}.  It contains
the profile, auxiliary, Green-factor, representer, and exact AU.1 output records
used by this theorem.  Its SHA--256 digest is
\begin{center}
\small\nolinkurl{1d5fe53786cc280006d7b1092d360556d4d8d8684e5ae3356ce8cd6d084e72a9}.
\end{center}
Equations~\eqref{eq:model-newton-bounds}--\eqref{eq:model-inertia-bound} are
conservative readable consequences of that archive; the self-map and
contraction decisions themselves are made before decimal formatting.  A
historical AU.1-only extraction is not a release dependency because its
embedded source snapshot predates the current audit driver.

\begin{corollary}[Local analytic parameter branch]
\label{cor:local-parameter-branch}
There is an open neighborhood
\[
 U\subset\{(m^2,\ell,R):m^2>0,\ \ell>0,\ \ell R^2<1\}
\]
of $(1,1/400,4)$ and a real-analytic family of regular solutions
$F_{m^2,\ell,R}$ of the corresponding Dirichlet problem on $[0,R]$.
After pullback to a fixed interval, this family is locally unique near
$F_*$.  On a possibly smaller neighborhood, every member is strictly
decreasing, has nonzero negative derivative at the wall, and has finite
positive rotor inertia.
\end{corollary}

\begin{proof}
Set $\theta=(m^2,\ell,R)$ and $\theta_0=(1,1/400,4)$.  Pull $[a,R]$
back to $I=[a,4]$ by
\[
 s_\theta=\frac{R-a}{4-a},
 \qquad x_\theta(\xi)=a+s_\theta(\xi-a).
\]
Let $\cY=L^\infty(I)\times\mathbb R$ and use the base graph coordinates
$X=(w,\beta)\in\cX$ from \eqref{eq:model-chart}.  Define the
parameter-dependent endpoint lift on the fixed interval by
\begin{equation}
 \widetilde F_{X,\theta}
 =\overline F+w+\beta H
  +\chi_a\left[
    \Phi_\theta(\overline b+\beta)-\Phi_{\theta_0}(\overline b)
    -\Phi'_{\theta_0}(\overline b)\beta\right].
 \label{eq:parameter-endpoint-chart}
\end{equation}
Because $H(a)=\Phi'_{\theta_0}(\overline b)$ and all three correction
directions vanish at $\xi=4$, this gives
\(\widetilde F_{X,\theta}(a)=\Phi_\theta(\overline b+\beta)\) and
\(\widetilde F_{X,\theta}(4)=0\) exactly.  At $\theta=\theta_0$ it reduces
to \eqref{eq:model-chart}.

For clarity, write the pulled-back coefficients as
\begin{align*}
 \widetilde{\mathsf p}_\theta(\xi,F)
 &=\bigl(1-\ell x_\theta(\xi)^2\bigr)
   \bigl(x_\theta(\xi)^2+8\sin^2F\bigr),\\
 \widetilde W_\theta(\xi,F)
 &=\sin^2F+\frac{2\sin^4F}{x_\theta(\xi)^2}
   +m^2x_\theta(\xi)^2(1-\cos F).
\end{align*}
The transformed interior operator is therefore
\begin{align}
 \mathcal G^{\rm pb}_\theta(\widetilde F)
 ={}&-s_\theta^{-2}\partial_\xi
       \left(\widetilde{\mathsf p}_\theta(\xi,\widetilde F)
             \partial_\xi\widetilde F\right)
 \notag\\
 &+\frac12s_\theta^{-2}
       \partial_F\widetilde{\mathsf p}_\theta(\xi,\widetilde F)
       (\partial_\xi\widetilde F)^2
   +\partial_F\widetilde W_\theta(\xi,\widetilde F).
 \label{eq:parameter-pulled-back-operator}
\end{align}
Thus the interior residual and the physical derivative-matching condition
define
\begin{equation}
 \mathscr R:\mathcal O_{\cX}\times\Theta\longrightarrow\cY,
 \qquad
 \mathscr R(X,\theta)=
 \left(\mathcal G^{\rm pb}_\theta(\widetilde F_{X,\theta}),
       s_\theta^{-1}\partial_\xi\widetilde F_{X,\theta}(a)
       -\Gamma_\theta(\overline b+\beta)\right).
 \label{eq:parameter-residual-map}
\end{equation}
The positive-radius coefficients make \eqref{eq:parameter-residual-map} real
analytic in $(X,\theta)$ while $\ell R^2<1$.  The regular-origin Volterra map is
jointly analytic in its shooting slope and $\theta$ and remains a strict
contraction after shrinking $\Theta$.  Its fixed point is therefore analytic,
so the cutoff maps $\Phi_\theta$ and $\Gamma_\theta$ are jointly analytic as
used in \eqref{eq:parameter-residual-map}.

At the exact zero $X_*$, the certificate gives
\[
 \norm{I-BD_X\mathscr R(X_*,\theta_0)}<0.31501<1.
\]
The reference block inverse $B$ is an isomorphism, so the Neumann series makes
$D_X\mathscr R(X_*,\theta_0)$ invertible.  The analytic implicit-function theorem now
gives the asserted local branch and local uniqueness.  The certificate bounds
$F_*'$ strictly below zero on every origin and positive-radius cell, including
the wall.  These inequalities persist under sufficiently small $C^1$
perturbations.  Positivity and finiteness of the inertia then persist by
continuity of its integrand on the compact pulled-back interval.
\end{proof}

\paragraph{Claim boundary.}
Theorem~\ref{thm:main} is an existence and local-uniqueness theorem for one
prescribed parameter point and one explicitly certified augmented
neighborhood.  Corollary~\ref{cor:local-parameter-branch} proves existence on
an open parameter neighborhood but does not give its numerical width or one
uniform Newton ball.  Neither statement is a global uniqueness theorem.  At
the certified point the geometry is fixed:
the result does not solve the Einstein--Skyrme equations or bound
gravitational backreaction.  The wall is an imposed ideal Dirichlet boundary,
and the theorem does not provide a support action, membrane dynamics, junction
conditions, finite-stiffness response, or wall inertia.  Nor does static existence establish
linear or nonlinear dynamical stability, an off-center solution, an
open-system approximation, or a de Sitter observer-algebra statement.
